\documentclass[../../main.tex]{subfiles}
\graphicspath{{../../image/}{../../figure/}}
\begin{document}

\section{注记}\label{section:ch02-注记}

连续物理这门历史悠久的学科，随着人们对非线性响应材料行为兴趣的兴起而正在经历复兴。Truesdell 与 Toupin [1] 以及 Truesdell 与 Noll [1] 的百科全书式著作包含可靠的历史信息以及大量参考文献书目，可作为追踪本学科从 18 世纪开端到 20 世纪 60 年代中期发展的优秀指南。Gurtin [1] 的教材对该领域提供了清晰、初等的介绍。更深入的处理并附带丰富文献可参见 Silhavy [1] 的书。Müller [2] 的教材是从现代连续物理视角对热力学的精彩阐述。其他强调弹性理论的优秀来源有 Ciarlet [1]、Hanyga [1]、Marsden 与 Hughes [1] 以及 Wang 与 Truesdell [1] 的著作。Antman [3] 的专著包含大量关于弹性弦、杆、壳和三维体理论的材料，重点放在支配平衡律的定性分析上。

运动的参考描述由 Euler 构想，并最终被命名为 Lagrange 描述，以凸显与 Lagrange 分析动力学表述的类比。另一方面，空间（即 Euler）描述由 Euler 有效使用，而最初由 Daniel Bernoulli 和 D'Alembert 引入。

关于连续介质物理平衡律场方程的参考（Lagrange）与空间（Euler）描述的等价性，见 Dafermos [17] 和 Wagner [2,3]。一个值得知道的问题是：这一等价性在比 Lipschitz 连续性更一般的运动假设下是否仍成立。例如，当介质是热弹性气体时，自然允许物体放置中存在真空区域。在这样的区域中密度为零，比容（变形梯度的行列式）变为无穷大。关于从 Euler 坐标过渡到 Lagrange 坐标时方程得到简化的其他例子，见 Peng [2]。

运动学平衡律\eqref{align:ch02-2-2-15} 和 \eqref{align:ch02-2-2-16} 最早由 Qin [1] 在光滑运动的框架中通过直接计算导出。有趣的是，正如我们在此所见，当运动仅仅 Lipschitz 连续时它们仍然成立，事实上 Demoulini、Stuart 和 Tzavaras [2] 已在更弱的假设下证明了这一点。与零 Lagrange 量的联系最早在该文献中指出。关于零 Lagrange 量的详细处理，见 Ball、Currie 和 Olver [1]。关于运动学平衡律的微分几何解释，见 Wagner [3,4]。

此处考虑的平衡律场方程最初由 Euler [1,2] 为质量、Cauchy [3,4] 为线动量和角动量、Kirchhoff [1] 为能量所导出。Clausius--Duhem 不等式由 Clausius [1] 对公绝热情形公设；熵通量项由 Duhem [1] 引入，熵产生项由 Truesdell 与 Toupin [1] 添加。更一般的熵不等式最早由 Müller [1] 考虑。

利用标架无差异和材料对称性来约化本构方程的做法，起源于 Cauchy [4] 和 Poisson [2] 的工作。在随后的一个世纪中，这一纲领由许多作者（大部分正确，但偶有错误）针对大量特殊本构方程加以实施。特别地，Cosserat 夫妇 [1]、Rivlin 与 Ericksen [1] 以及其他学者在 20 世纪 40 年代和 50 年代的工作促进了概念的澄清。材料标架无差异原理和对称群的定义最终由 Noll [1] 以充分的一般性和数学严密性公设。

本构方程应加以约化，使得任何平衡了质量、动量和能量的光滑热力学过程都自动满足 Clausius--Duhem 不等式，这一公设最早由 Coleman 与 Noll [1] 作为一般原理提出。此处给出的例子改编自 Coleman 与 Noll [1]（热弹性情形）以及 Coleman 与 Mizel [1]（热粘弹性情形）。

van der Waals 在其博士论文（1873）中引入了以其名字命名的状态方程，以说明气体分子体积和分子间力。Gallavotti [1] 从统计物理的视角讨论了其解释。van der Waals 气体多年来一直作为相变的简单模型。

正如我们在历史导论中所见，Chaplygin 气体的特殊性质最早由 Earnshaw [1] 注意到。Chaplygin [1] 的经典贡献在 von Mises [1] 的教材中得到阐述。Chaplygin 气体目前在宇宙学中找到了新的应用。关于该气体状态方程在微分几何中的一个出人意料的应用，见\cref{chapter:ch18-多维定常与自相似解}。


Coleman 与 Gurtin [1] 发展了带内部状态变量的热粘弹性材料一般理论，其中\cref{section:ch02-松弛} 给出的例子是一个特例。这类本构关系最早由 Maxwell [1] 考虑。关于气体动力学中松弛现象的详细讨论，见 Vincenti 与 Kruger [1] 的书。

\end{document}
